\documentclass[10pt,a4paper]{amsart}
\usepackage[margin=19mm]{geometry}
\usepackage{amsmath,amssymb,mathtools}
\usepackage[scr=rsfs,cal=euler]{mathalfa}
\usepackage{xfrac}
\usepackage[unicode,hidelinks,bookmarks=false]{hyperref}
\newcommand{\ie}{i.e.\ }
\newcommand{\cf}{cf.\ }
\newcommand{\resp}{respectively\ }
\newcommand{\Subsection}{\subsection}
\providecommand{\nobreakdash}{\nobreak\hbox{-}}
\setlength{\emergencystretch}{2em}
\multlinegap0pt
\usepackage{bm}
\usepackage{bbm}
\usepackage{dsfont}
\usepackage{xspace}
\usepackage{cancel}
\usepackage{mathtools}
\usepackage{amsthm}
\makeatletter
\@ifundefined{lemma}{\newtheorem{lemma}{Lemma}}{}
\makeatother
\newcommand{\As}{\mathscr{A}}
\newcommand{\Bf}{\mathfrak{B}}
\newcommand{\Cb}{\mathbb{C}}
\newcommand{\alg}{{\rm alg}}
\newcommand{\bra}[1]{\left< #1 \right|}
\newcommand{\ket}[1]{\left| #1 \right>}
\newcommand{\one}{\mathbb{1}}
\title[Quantum model reduction for continuous-time quantum filters]{Quantum model reduction for continuous-time quantum filters}
\author{Tommaso Grigoletto, Cl\'ement Pellegrini, Francesco Ticozzi}
\date{Source: arXiv:2501.13885v3 (CC BY 4.0)}
\begin{document}
\maketitle

\noindent\emph{Source and licence.} Quantum model reduction for continuous-time quantum filters. Version arXiv:2501.13885v3 (CC BY 4.0), distributed under the Creative Commons Attribution 4.0 International licence, as recorded in the arXiv licence metadata for this exact version and shown on the abstract page https://arxiv.org/abs/2501.13885v3. Full rights notice: licenses/arxiv-2501.13885-CC-BY-4.0.txt.\par\smallskip

\noindent\textit{Editorial scope.} Counted unit: the Lemma whose source label is \texttt{None} in the article (source0.tex lines 1118--1140, source label \texttt{None}), delivered together with the article's own complete original proof of it (source0.tex lines 1141--1177, one \texttt{proof} environment of 37 lines). No mathematics is written, repaired or re-derived: statement and proof are the article's own text line for line. Required context retained uncounted: none. Every definition, display and label of those lines is retained unchanged. Omitted from this package and not redistributed: 1--1117, 1178--1498. Nothing inside a retained excerpt is omitted or altered: every retained body is delivered exactly as the article's lines are written, so no omission receipt of any kind accompanies this package. Citations: the retained excerpts carry no citation command at all, so no bibliography entry of the article is transcribed and no reference is redistributed. Reference adaptation: none was needed and none was made; every reference inside the delivered unit resolves inside it, so no unresolved reference or citation is produced. Printed slips kept literal: no printed word, symbol, hypothesis or number of the counted statement or of its proof was altered. Layout and class: the article's own class is \texttt{imsart} with packages including natbib, hyperref, graphicx, xcolor, cancel, amsfonts, amsmath, amssymb, babel, ulem, amsthm, xfrac, algorithm2e, subcaption; the wrapper is the lane's pinned \texttt{amsart} editorial wrapper at 10pt on A4, which restates the theorem-like environments the retained text needs and adds the article's own macro definitions that the retained text needs (\texttt{\textbackslash As}, \texttt{\textbackslash Bf}, \texttt{\textbackslash Cb}, \texttt{\textbackslash alg}, \texttt{\textbackslash bra}, \texttt{\textbackslash ket}, \texttt{\textbackslash one}), with extra wrapper packages (bm, bbm, dsfont, xspace, cancel, mathtools). The delivered numbering is the unit's own. Assets: no retained span contains a figure environment, a graphics-inclusion command, a PDF-page inclusion, a table or a TikZ picture; no image file is copied into this package, the package declares no asset dependency and no image file is redistributed with it. Licence: version arXiv:2501.13885v3 of this article is distributed under the Creative Commons Attribution 4.0 International licence, as recorded in the arXiv licence metadata for this exact version and shown on the abstract page https://arxiv.org/abs/2501.13885v3. A full rights notice is delivered at licenses/arxiv-2501.13885-CC-BY-4.0.txt. Characters: every retained excerpt is pure ASCII, so the delivered engine, XeLaTeX with its default Latin Modern text font, reports no missing character.
\begin{lemma}
    Let us define the permutation matrix \[P \equiv \one_4 + \sigma_x\otimes\one_2 + \one_2\otimes\sigma_z  -\sigma_x\otimes\sigma_z = \begin{bmatrix}
        1&0&0&0\\0&0&0&1\\0&0&1&0\\0&1&0&0
    \end{bmatrix} \in\Cb^{4\times4}\] and define, for any $N\geq1$, the unitary operator \[U_N = \begin{cases}
        \one_2 & \text{if }N=1\\
        (P\otimes\one_{2^{N-2}})(\one_2\otimes U_{N-1}) &\text{if } N\geq2
    \end{cases}.\]

    Then for all $N$: 
    \begin{enumerate}
        \item \[U_N \sigma_z^{(j)} U_N^* = \begin{cases}
            \sigma_z^{(j)}\sigma_z^{(j+1)}&\text{if }j<N\\
            \sigma_z^{(N)}&\text{if }j=N
        \end{cases};\]
        \item \(U_N \sigma_x^{(j)}\sigma_x^{(j)} U_N^* = \sigma_x^{(j+1)}\);
        \item $U_N \As U_N^* = \Bf(\Cb^{2^{N-1}}) \bigoplus \Bf(\Cb^{2^{N-1}})$;
        \item 
        \[U_N \sigma_-^{(j)} U_N^* =\begin{cases}
            \frac{1}{2}\left[\sigma_x^{(0)}\sigma_x^{(1)}\dots\sigma_x^{(j-1)}\sigma_x^{(j)} -i \sigma_x^{(0)}\sigma_x^{(1)}\dots\sigma_x^{(j-1)}\sigma_y^{(j)} \sigma_z^{(j+1)} \right]&\text{if }j < N\\
            \sigma_x^{(0)}\sigma_x^{(1)}\dots\sigma_x^{(N-1)}\sigma_-^{(N)}&\text{if }j=N
        \end{cases}. \]
    \end{enumerate}
\end{lemma}

\begin{proof}
Note that the fact that $U_N$ is unitary can easily proven by induction.
    The first claim of this Lemma is also proven by induction. 
    We start by proving the case $N=2$. Simple calculations show that 
    \begin{align*}
        U_2\sigma_z^{(1,2)} U_2^* &= P (\sigma_z\otimes\one_2) P^* = \sigma_z\otimes\sigma_z\\
        U_2\sigma_z^{(2,2)} U_2^* &= P (\one_2\otimes\sigma_z) P^* = \one_2\otimes\sigma_z.
    \end{align*}
    Now assume that $U_{N-1}\sigma_z^{(j,N-1)}U_{N-1}^* = \sigma_z^{(j,N-1)}\sigma_z^{(j,N-1)}$ for $j<N-1$ and $U_{N-1}\sigma_z^{(N-1,N-1)}U_{N-1}^* = \sigma_z^{(N-1,N-1)}$. Then,
    \begin{align*}
        U_N\sigma_z^{(1,N)} U_N^* &= (P\otimes\one_{2^{N-2}})(\one_2\otimes U_{N-1}) (\sigma_z\otimes \one_{2^{N-1}}) (\one_2\otimes U_{N-1}^*) (P^*\otimes\one_{2^{N-2}})\\
        &= P(\sigma_z\otimes\one_2)P^* \otimes\one_{2^{N-2}} = \sigma_z\otimes\sigma_z\otimes\one_{2^{N-2}} = \sigma_z^{(1,N)}\sigma_z^{(2,N)},
    \end{align*}
    and, for $j\geq2$
    \begin{align*}
        U_N\sigma_z^{(j,N)} U_N^* &= (P\otimes\one_{2^{N-2}})(\one_2\otimes U_{N-1}) (\one_2\otimes \sigma_z^{(j-1,N-1)}) (\one_2\otimes U_{N-1}^*) (P^*\otimes\one_{2^{N-2}})\\
        &= (P\otimes\one_{2^{N-2}})(\one_2\otimes \sigma_z^{(j-1,N-1)}\sigma_z^{(j,N-1)}) (P^*\otimes\one_{2^{N-2}}) = (\one_2\otimes \sigma_z^{(j-1,N-1)}\sigma_z^{(j,N-1)}) = \sigma_z^{(j,N)}\sigma_z^{(j+1,N)}
    \end{align*}
    concluding the proof of the third point. 

    The next point is also proven by induction. 
    We then prove the base case $N=2$. Simple calculations show that \(U_2\sigma_x^{(1,2)}\sigma_x^{(2,2)} U_2^* = P (\sigma_x\otimes\sigma_x) P^* = \one_2\otimes\sigma_x\) concluding the base case. Similar calculations show that \(P(\sigma_x\otimes\one_2)P^* = \sigma_x\otimes \one_2 \) and \(P(\one_2\otimes\sigma_x)P^* = \sigma_x\otimes \sigma_x \) which are useful in what comes next. We then assume that $U_{N-1} \sigma_x^{(j,N-1)}\sigma_x^{(j,N-1)} U_{N-1^*} = \sigma_x^{(j+1,N-1)}$ and consider
    \begin{align*}
        U_N \sigma_x^{(1,N)}\sigma_x^{(2,N)}U_N^* &= (P\otimes\one_{2^{N-2}})(\one_2\otimes U_{N-1}) (\sigma_x\otimes \sigma_x^{(2,N-1)}) (\one_2\otimes U_{N-1}^*) (P^*\otimes\one_{2^{N-2}})\\
        &= (P\otimes\one_{2^{N-2}}) (\sigma_x\otimes \sigma_x^{(1,N-1)}) (P^*\otimes\one_{2^{N-2}}) = (P\otimes\one_{2^{N-2}}) (\sigma_x\otimes \sigma_x \otimes \one_{2^{N-2}} ) (P^*\otimes\one_{2^{N-2}}) \\&= \one_2\otimes\sigma_x\otimes \one_{2^{N-2}} = \sigma_x^{(2,N)}
    \end{align*}
    where we used the fact that $U_N\sigma_x^{(j,N)}U_N^* = \sigma_x^{(j,N)}$ which can also be proved by induction, and, for $j>1$:
    \begin{align*}
        U_N \sigma_x^{(j,N)}\sigma_x^{(j+1,N)}U_N^* &= (P\otimes\one_{2^{N-2}})(\one_2\otimes U_{N-1}) (\one_2\otimes \sigma_x^{(j-1,N-1)}\sigma_x^{(j,N-1)}) (\one_2\otimes U_{N-1}^*) (P^*\otimes\one_{2^{N-2}})\\
        &= (P\otimes\one_{2^{N-2}})(\one_2\otimes \sigma_x^{(j,N-1)}) (P^*\otimes\one_{2^{N-2}}) = \sigma_x^{(j+1,N-1)}
    \end{align*}
    which concludes the proof of the second statement. 

    The third claim is a direct consequence of the first two claims as, by definition $\As = \alg(\{\sigma_z^{(j)}\}_{j=1}^N\cup\{\sigma_x^{(j)}\sigma_x^{(j+1)}\}_{j=1}^{N-1})$ and the fact that $U_N \sigma_z^{(j)} U_N^*$ and $U_N \sigma_x^{(j)}\sigma_x^{(j)} U_N^*$ act on the first qubit either as the identity operator or as $\sigma_z$, thus the off-diagonal blocks must be zero, e.g. $\bra{k}\otimes\one_{2^{N-1}} \sigma_x^{(j)} \ket{l}\otimes\one_{2^{N-1}} = 0$.

    The fourth claim is proven by induction as the first two claims and its proof is here omitted.
\end{proof}

\end{document}
